\documentclass[11pt,a4paper]{article}
\usepackage[utf8]{inputenc}
\usepackage[margin=2.5cm]{geometry}
\usepackage{amsmath}
\usepackage{amssymb}
\usepackage{graphicx}
\usepackage{hyperref}
\usepackage{enumitem}
\usepackage{booktabs}

\title{\textbf{Controller Optimization System for Manufacturing Process Chain}\\
\large Uncertainty Quantification and Policy Learning}
\author{AZIMUTH Project}
\date{\today}

\begin{document}

\maketitle

\begin{abstract}
This document describes the conceptual architecture and methodology of a two-phase controller optimization system for sequential manufacturing processes. The system combines uncertainty quantification with reinforcement learning-inspired policy optimization to improve manufacturing reliability. Phase 1 trains process-specific uncertainty predictors using Structural Causal Models (SCMs). Phase 2 optimizes policy generators that adapt process inputs to minimize uncertainty propagation while maintaining product quality.
\end{abstract}

\section{Introduction}

Modern manufacturing involves sequential processes where the output of one stage becomes the input to the next. In this work, we consider a three-stage chain:

\begin{enumerate}
    \item \textbf{Laser Processing}: Precision cutting or surface treatment
    \item \textbf{Plasma Treatment}: Surface modification or cleaning
    \item \textbf{Galvanic Deposition}: Metal coating application
\end{enumerate}

Each process introduces uncertainty due to environmental variations, material inconsistencies, and equipment tolerances. Our goal is to develop an adaptive controller that:
\begin{itemize}
    \item Learns to predict process uncertainties
    \item Generates optimal control inputs to minimize uncertainty propagation
    \item Maximizes final product reliability
\end{itemize}

\section{System Architecture}

The system operates in two distinct phases, following a curriculum learning approach:

\subsection{Phase 1: Uncertainty Predictor Training}

For each manufacturing process $p \in \{\text{laser, plasma, galvanic}\}$, we train an \textbf{Uncertainty Predictor} $\mathcal{U}_p$ that maps control inputs to output distributions.

\subsubsection{Input-Output Relationship}
Given control inputs $\mathbf{x}_p \in \mathbb{R}^{d_{in}}$, the uncertainty predictor outputs:
\begin{align}
\mathcal{U}_p(\mathbf{x}_p) = (\boldsymbol{\mu}_p, \boldsymbol{\sigma}^2_p)
\end{align}
where $\boldsymbol{\mu}_p \in \mathbb{R}^{d_{out}}$ is the mean output and $\boldsymbol{\sigma}^2_p \in \mathbb{R}^{d_{out}}$ is the variance (uncertainty).

\subsubsection{Training Data Generation}
Training data is generated using \textbf{Structural Causal Models (SCMs)} that simulate realistic manufacturing processes with:
\begin{itemize}
    \item Physical equations modeling process dynamics
    \item Stochastic noise terms representing environmental variations
    \item Causal relationships between variables
\end{itemize}

\subsubsection{Neural Network Architecture}
Each uncertainty predictor is a neural network with:
\begin{itemize}
    \item Shared hidden layers with batch normalization and dropout
    \item Dual output heads:
    \begin{itemize}
        \item Mean head: predicts expected output $\boldsymbol{\mu}_p$
        \item Variance head: predicts aleatoric uncertainty $\boldsymbol{\sigma}^2_p$ (constrained positive)
    \end{itemize}
\end{itemize}

\subsubsection{Loss Function}
The network is trained using \textbf{Negative Log-Likelihood (NLL)} loss:
\begin{align}
\mathcal{L}_{NLL} = \frac{1}{N} \sum_{i=1}^{N} \left[ \frac{1}{2}\log(\sigma_i^2) + \frac{(y_i - \mu_i)^2}{2\sigma_i^2} \right]
\end{align}
This loss encourages the model to:
\begin{itemize}
    \item Minimize prediction error $|y_i - \mu_i|$
    \item Calibrate uncertainty $\sigma_i^2$ appropriately (not too large, not too small)
\end{itemize}

\subsubsection{Training Output}
For each process, Phase 1 produces:
\begin{itemize}
    \item Trained uncertainty predictor weights (frozen for Phase 2)
    \item Input/output scalers for normalization
    \item Performance metrics (MSE, MAE, R$^2$, calibration ratio)
    \item PDF report with visualizations
\end{itemize}

\subsection{Phase 2: Policy Generator Training}

In this phase, we train \textbf{Policy Generators} $\pi_p$ that learn to produce optimal control inputs for each process, given the current system state.

\subsubsection{Process Chain Architecture}
The complete process chain consists of:
\begin{enumerate}
    \item \textbf{Frozen Uncertainty Predictors} $\mathcal{U}_{\text{laser}}, \mathcal{U}_{\text{plasma}}, \mathcal{U}_{\text{galvanic}}$ from Phase 1
    \item \textbf{Trainable Policy Generators}:
    \begin{itemize}
        \item $\pi_{\text{plasma}}$: generates plasma inputs from laser outputs
        \item $\pi_{\text{galvanic}}$: generates galvanic inputs from plasma outputs
    \end{itemize}
\end{enumerate}

Note: The laser process receives external inputs, so no policy generator is needed for it.

\subsubsection{Forward Pass Through Process Chain}
\begin{enumerate}
    \item External inputs $\mathbf{x}_{\text{laser}}$ enter laser process
    \item Laser uncertainty predictor: $(\boldsymbol{\mu}_{\text{laser}}, \boldsymbol{\sigma}^2_{\text{laser}}) = \mathcal{U}_{\text{laser}}(\mathbf{x}_{\text{laser}})$
    \item Plasma policy generator: $\mathbf{x}_{\text{plasma}} = \pi_{\text{plasma}}(\boldsymbol{\mu}_{\text{laser}})$
    \item Plasma uncertainty predictor: $(\boldsymbol{\mu}_{\text{plasma}}, \boldsymbol{\sigma}^2_{\text{plasma}}) = \mathcal{U}_{\text{plasma}}(\mathbf{x}_{\text{plasma}})$
    \item Galvanic policy generator: $\mathbf{x}_{\text{galvanic}} = \pi_{\text{galvanic}}(\boldsymbol{\mu}_{\text{plasma}})$
    \item Galvanic uncertainty predictor: $(\boldsymbol{\mu}_{\text{galvanic}}, \boldsymbol{\sigma}^2_{\text{galvanic}}) = \mathcal{U}_{\text{galvanic}}(\mathbf{x}_{\text{galvanic}})$
\end{enumerate}

The complete trajectory is defined as:
\begin{align}
\mathbf{a} = \{(\mathbf{x}_p, \boldsymbol{\mu}_p, \boldsymbol{\sigma}^2_p) : p \in \{\text{laser, plasma, galvanic}\}\}
\end{align}

\subsubsection{Training Trajectories}
We generate three reference trajectories for comparison:

\begin{enumerate}
    \item \textbf{Target Trajectory} $\mathbf{a}^*$: Generated using SCMs with near-zero noise ($\epsilon = 10^{-8}$). Represents the ideal, deterministic behavior.

    \item \textbf{Baseline Trajectory} $\mathbf{a}'$: Generated using SCMs with normal process noise. Represents uncontrolled operation without policy generators.

    \item \textbf{Actual Trajectory} $\mathbf{a}$: Generated by the process chain with policy generators during training. Represents controlled operation.
\end{enumerate}

\section{Surrogate Reliability Model}

A key component is the \textbf{Surrogate Model} that evaluates trajectory quality.

\subsection{Reliability Function}
The surrogate computes a reliability score $F(\mathbf{a})$ for any trajectory $\mathbf{a}$. Higher $F$ indicates better quality.

\subsubsection{Sampling-Based Evaluation}
For each process output distribution $(\boldsymbol{\mu}_p, \boldsymbol{\sigma}^2_p)$, we sample:
\begin{align}
\mathbf{y}_p \sim \mathcal{N}(\boldsymbol{\mu}_p, \boldsymbol{\sigma}^2_p)
\end{align}
using the reparameterization trick:
\begin{align}
\mathbf{y}_p = \boldsymbol{\mu}_p + \boldsymbol{\epsilon} \odot \sqrt{\boldsymbol{\sigma}^2_p}, \quad \boldsymbol{\epsilon} \sim \mathcal{N}(0, I)
\end{align}
This maintains differentiability for backpropagation.

\subsubsection{Quality Metrics}
For each process, we compute a quality score based on proximity to optimal values:
\begin{align}
q_{\text{laser}} &= \exp\left(-\frac{(y_{\text{laser}} - 0.5)^2}{0.1}\right) \\
q_{\text{plasma}} &= \exp\left(-\frac{(y_{\text{plasma}} - 5.0)^2}{2.0}\right) \\
q_{\text{galvanic}} &= \exp\left(-\frac{(y_{\text{galvanic}} - 10.0)^2}{4.0}\right)
\end{align}

where the optimal values (0.5, 5.0, 10.0) and variances (0.1, 2.0, 4.0) are domain-specific parameters.

\subsubsection{Weighted Combination}
The final reliability score is a weighted combination emphasizing the final product:
\begin{align}
F(\mathbf{a}) = 0.2 \cdot q_{\text{laser}} + 0.3 \cdot q_{\text{plasma}} + 0.5 \cdot q_{\text{galvanic}}
\end{align}

\subsection{Target Reliability}
We compute $F^* = F(\mathbf{a}^*)$ once at initialization. Since $\mathbf{a}^*$ has near-zero uncertainty and optimal parameter values, $F^*$ represents the best achievable reliability.

\section{Policy Optimization}

\subsection{Loss Function}
The policy generators are trained using a composite loss:
\begin{align}
\mathcal{L}_{\text{total}} = \mathcal{L}_{\text{reliability}} + \lambda \cdot \mathcal{L}_{\text{BC}}
\end{align}

\subsubsection{Reliability Loss}
Encourages the actual trajectory to achieve high reliability:
\begin{align}
\mathcal{L}_{\text{reliability}} = (F(\mathbf{a}) - F^*)^2
\end{align}
Minimizing this pushes $F(\mathbf{a})$ toward the optimal $F^*$.

\subsubsection{Behavior Cloning Loss}
Regularizes the policy to stay close to target inputs:
\begin{align}
\mathcal{L}_{\text{BC}} = \sum_{p} \|\mathbf{x}_p - \mathbf{x}_p^*\|^2
\end{align}
This prevents the policy from exploring unrealistic input regions.

\subsection{Training Algorithm}
\begin{enumerate}
    \item Sample a batch of initial laser inputs
    \item Forward pass through process chain to obtain $\mathbf{a}$
    \item Compute $F(\mathbf{a})$ using surrogate model
    \item Compute loss $\mathcal{L}_{\text{total}}$
    \item Backpropagate gradients \textbf{only through policy generators}
    \item Update policy parameters using optimizer (Adam)
\end{enumerate}

Note: Uncertainty predictor parameters remain frozen throughout Phase 2.

\section{Evaluation Metrics}

\subsection{Reliability Comparison}
We compare three reliability scores:
\begin{itemize}
    \item $F^*$: Target reliability (optimal, near-zero noise)
    \item $F'$: Baseline reliability (uncontrolled, normal noise)
    \item $F$: Actual reliability (controlled, with policy generators)
\end{itemize}

\textbf{Success criteria}:
\begin{itemize}
    \item $F > F'$: Controller improves over baseline
    \item $F \approx F^*$: Controller approaches optimal performance
\end{itemize}

\subsection{Trajectory Distance Metrics}
For each pair of trajectories, we compute:
\begin{align}
d(\mathbf{a}_1, \mathbf{a}_2) = \sum_p \left( \|\mathbf{x}_{p,1} - \mathbf{x}_{p,2}\|^2 + \|\boldsymbol{\mu}_{p,1} - \boldsymbol{\mu}_{p,2}\|^2 \right)
\end{align}

Key comparisons:
\begin{itemize}
    \item $d(\mathbf{a}, \mathbf{a}^*)$: How close is the controlled trajectory to optimal?
    \item $d(\mathbf{a}', \mathbf{a}^*)$: How close is the uncontrolled trajectory to optimal?
    \item $d(\mathbf{a}, \mathbf{a}')$: How different is controlled vs. uncontrolled?
\end{itemize}

\subsection{Process-Wise Analysis}
For each individual process, we report:
\begin{itemize}
    \item Input deviation: $\|\mathbf{x}_p - \mathbf{x}_p^*\|$
    \item Output deviation: $\|\boldsymbol{\mu}_p - \boldsymbol{\mu}_p^*\|$
    \item Uncertainty magnitude: $\|\boldsymbol{\sigma}_p^2\|$
\end{itemize}

\section{Results Interpretation}

\subsection{Training Convergence}
Monitor the training history for:
\begin{itemize}
    \item Total loss decreasing over epochs
    \item Reliability loss: $(F - F^*)^2$ approaching zero
    \item Behavior cloning loss: stabilizing at reasonable level
    \item $F$ (actual reliability) increasing toward $F^*$
\end{itemize}

\subsection{Performance Indicators}

\textbf{Improvement over baseline}:
\begin{align}
\text{Improvement} = \frac{F - F'}{F'} \times 100\%
\end{align}
Positive values indicate successful controller learning.

\textbf{Gap from optimal}:
\begin{align}
\text{Gap} = \frac{F^* - F}{F^*} \times 100\%
\end{align}
Lower values indicate better approximation of optimal behavior.

\subsection{Visualization Outputs}
The system generates:
\begin{enumerate}
    \item \textbf{Training History Plot}: Loss evolution over epochs
    \item \textbf{Trajectory Comparison}: Side-by-side visualization of $\mathbf{a}^*$, $\mathbf{a}'$, $\mathbf{a}$
    \item \textbf{Reliability Comparison}: Bar chart of $F^*$, $F'$, $F$
    \item \textbf{Process Improvements}: Per-process metric breakdown
    \item \textbf{PDF Report}: Comprehensive summary with all metrics and plots
\end{enumerate}

\section{Key Design Principles}

\subsection{Modularity}
\begin{itemize}
    \item Each process has its own uncertainty predictor (separate weights)
    \item Policy generators are process-specific
    \item Easy to add/remove processes from the chain
\end{itemize}

\subsection{Transfer Learning}
\begin{itemize}
    \item Phase 1 uncertainty predictors are pre-trained on SCM data
    \item Phase 2 leverages frozen predictors, focusing only on policy learning
    \item Reduces sample complexity and training time
\end{itemize}

\subsection{Uncertainty Awareness}
\begin{itemize}
    \item Policies receive uncertainty information from previous stages
    \item Can adapt inputs to compensate for upstream variability
    \item Minimizes uncertainty propagation through the chain
\end{itemize}

\subsection{Differentiable Simulation}
\begin{itemize}
    \item Entire process chain is differentiable
    \item Enables gradient-based policy optimization
    \item Reparameterization trick maintains differentiability through stochastic sampling
\end{itemize}

\section{Future Extensions}

\subsection{Real Process Integration}
Replace SCM-generated data with real sensor measurements from manufacturing equipment.

\subsection{Online Adaptation}
Implement continual learning to adapt policies as process conditions drift over time.

\subsection{Multi-Objective Optimization}
Extend to optimize multiple objectives: reliability, energy efficiency, throughput, cost.

\subsection{ProT Surrogate Model}
Replace the placeholder surrogate with a trained transformer model (ProT) that learns complex quality functions from data.

\subsection{Constraint Handling}
Add explicit constraints on inputs/outputs (e.g., safety limits, material compatibility).

\section{Conclusion}

This two-phase controller optimization system demonstrates a principled approach to managing uncertainty in sequential manufacturing processes. By separating uncertainty quantification (Phase 1) from policy learning (Phase 2), the system achieves:

\begin{itemize}
    \item Reliable uncertainty estimates calibrated to process physics
    \item Data-efficient policy learning leveraging pre-trained models
    \item Interpretable results through comprehensive metrics and visualizations
    \item Scalable architecture adaptable to various manufacturing chains
\end{itemize}

The combination of neural uncertainty quantification and gradient-based policy optimization provides a foundation for intelligent, adaptive manufacturing control systems.

\end{document}
